\documentclass[11pt]{article}
\usepackage[margin=1in]{geometry}
\usepackage{graphicx}
\usepackage{booktabs}
\usepackage{array}
\usepackage{float}
\usepackage{enumitem}
\usepackage{xcolor}

\setlength{\parskip}{0.5em}
\setlength{\parindent}{0pt}
\newcommand{\code}[1]{\texttt{#1}}

% The figures are in figures/ next to this file. If that folder is missing,
% the document still compiles and shows a note in place of each figure; the
% submitted project_report.pdf contains all of them.
\newcommand{\figurefile}[2]{%
  \IfFileExists{#2}{\includegraphics[width=#1]{#2}}{%
    \fbox{\parbox[c][3cm][c]{\dimexpr #1-2\fboxsep-2\fboxrule\relax}{\centering\small
      Figure file \texttt{\detokenize{#2}} is not included here;\\
      see \texttt{project\_report.pdf} for this figure.}}}}

\title{CSC258 Project Report: Magical Drop}
\author{Xiangyu Zhou (1012571708) \and Belina Zhu (TODO student number)}
\date{}

\begin{document}
\maketitle

\section{Summary}

\code{MagicalDrop.asm} is a MIPS assembly version of Magical Drop for the
Bitmap Display and the keyboard of Saturn or MARS. It reaches Milestones 1 to
5. For Milestones 4 and 5 it uses combination (d) of Milestone 5: two easy
features and three hard features (Table~\ref{tab:features}, numbered as in the
project handout). The same five features also satisfy combination (c) of
Milestone 4.

\begin{table}[H]
\centering
\begin{tabular}{@{}clp{9.5cm}@{}}
\toprule
 & No. & Feature\\
\midrule
Easy & 1 & Row advancement on a timer: the first timer row comes after 500 frames\\
Easy & 2 & Row advancement gets faster: 50 frames less after every new row, down to 150\\
\midrule
Hard & 2 & Flood fill (horizontal and vertical) after the vertical match-3 trigger\\
Hard & 3 & Gravity: unsupported bubbles float upwards\\
Hard & 6 & A flashing dotted aiming line in the character's column\\
\bottomrule
\end{tabular}
\caption{Features implemented for Milestones 4 and 5.}
\label{tab:features}
\end{table}

Hard features 2 and 3 are steps of the matching process of
Section~\ref{sec:m3}; hard feature 6 is part of the drawing code.

\section{Bitmap Display Configuration}

\begin{center}
\begin{tabular}{@{}ll@{}}
\toprule
Unit width in pixels & 8\\
Unit height in pixels & 8\\
Display width in pixels & 256\\
Display height in pixels & 512\\
Base address for display & \code{0x10008000} (\code{\$gp})\\
\bottomrule
\end{tabular}
\end{center}

The display is therefore $32 \times 64$ units. Unit $(x, y)$ is the word at
address \code{0x10008000} $+\ 4\,(32y + x)$, so moving one unit to the right
adds 4 bytes and moving one unit down adds 128 bytes. The display uses
$2048$ words (\code{0x10008000} to \code{0x10009FFF}), so it ends before the
\code{.data} section at \code{0x10010000} and no padding variable is needed.

\section{Memory Layout Plan}

\subsection{What is drawn where}

Figure~\ref{fig:layout} is the plan of the screen. One square of the plan is
one unit of the display. A bubble is $3 \times 3$ units and the cells of the
playing field are 4 units apart, which leaves one black unit between
neighbouring bubbles. The cell in row $r$ and column $c$ has its top-left
unit at
\[
x = 2 + 4c, \qquad y = 2 + 4r \qquad (r = 0,\dots,13,\quad c = 0,\dots,6).
\]

\begin{figure}[H]
\centering
\figurefile{0.9\textwidth}{figures/layout.png}
\caption{Plan of the $32 \times 64$ unit display.}
\label{fig:layout}
\end{figure}

\subsection{The playing field}

The playing field is the $7 \times 14$ grid required by the handout. It is
stored in the array \code{grid} as 98 words in row major order, separate from
the character and the bubbles it holds. Each word is \code{0} for an empty
cell or the colour of the bubble: \code{1} red, \code{2} yellow, \code{3}
green, \code{4} blue. Cell $(r, c)$ has index $7r + c$ and is stored at
\code{grid} $+\ 4\,(7r + c)$. Table~\ref{tab:grid} shows the index of every
cell.

Row 0 is the top row. Bubbles hang from the top, so in every column the
bubbles occupy rows $0, 1, 2, \dots$ without gaps. Row 13 is the character's
row: bubbles can only get there when a new row pushes the field down, and
that ends the game.

\begin{table}[H]
\centering
\begin{tabular}{@{}r|*{7}{>{\centering\arraybackslash}p{0.9cm}}|l@{}}
 & $c=0$ & 1 & 2 & 3 & 4 & 5 & 6 & \\
\midrule
$r=0$ & 0 & 1 & 2 & 3 & 4 & 5 & 6 & top row (new rows arrive here)\\
1 & 7 & 8 & 9 & 10 & 11 & 12 & 13 & \\
2 & 14 & 15 & 16 & 17 & 18 & 19 & 20 & \\
3 & 21 & 22 & 23 & 24 & 25 & 26 & 27 & \\
4 & 28 & 29 & 30 & 31 & 32 & 33 & 34 & last row filled at the start\\
5 & 35 & 36 & 37 & 38 & 39 & 40 & 41 & \\
6 & 42 & 43 & 44 & 45 & 46 & 47 & 48 & \\
7 & 49 & 50 & 51 & 52 & 53 & 54 & 55 & \\
8 & 56 & 57 & 58 & 59 & 60 & 61 & 62 & \\
9 & 63 & 64 & 65 & 66 & 67 & 68 & 69 & \\
10 & 70 & 71 & 72 & 73 & 74 & 75 & 76 & \\
11 & 77 & 78 & 79 & 80 & 81 & 82 & 83 & \\
12 & 84 & 85 & 86 & 87 & 88 & 89 & 90 & lowest row a bubble may rest in\\
\midrule
13 & 91 & 92 & 93 & 94 & 95 & 96 & 97 & the character's row\\
\end{tabular}
\caption{Index of each cell in \code{grid} (and in \code{mark}). The word of
cell $i$ is at \code{grid} $+\ 4i$.}
\label{tab:grid}
\end{table}

\subsection{Variables}

Table~\ref{tab:memory} lists everything in the \code{.data} section in the
order in which it is laid out. All game state lives in memory; registers only
hold temporary values.

\begin{table}[H]
\centering
\small
\begin{tabular}{@{}lllp{7.2cm}@{}}
\toprule
Address & Label & Size & Contents\\
\midrule
\code{0x10010000} & \code{display\_address} & 1 word & \code{0x10008000}\\
\code{0x10010004} & \code{keyboard\_address} & 1 word & \code{0xffff0000}\\
\code{0x10010008} & \code{grid} & 98 words & the playing field (Table~\ref{tab:grid})\\
\code{0x10010190} & \code{mark} & 98 words & 1 for every cell of the match being removed\\
\code{0x10010318} & \code{ff\_stack} & 98 words & work stack of the flood fill (cell indices)\\
\code{0x100104a0} & \code{ff\_count} & 1 word & number of entries on \code{ff\_stack}\\
\code{0x100104a4} & \code{player\_col} & 1 word & column of the character, 0 to 6 (starts at 3)\\
\code{0x100104a8} & \code{held\_colour} & 1 word & colour of the held bubbles, 0 if none\\
\code{0x100104ac} & \code{held\_count} & 1 word & number of held bubbles\\
\code{0x100104b0} & \code{row\_timer} & 1 word & frames since the last new row\\
\code{0x100104b4} & \code{key\_count} & 1 word & keystrokes (\code{a}, \code{d}, \code{s}, \code{w}) since the last new row\\
\code{0x100104b8} & \code{row\_interval} & 1 word & frames until the next timer row: 500 at the start, 50 less after every new row, at least 150\\
\code{0x100104bc} & \code{frame\_count} & 1 word & frames drawn so far (makes the aiming line flash)\\
\midrule
\code{0x100104c0} & \code{palette} & 15 words & 5 entries of 3 shades (dark, base, light)\\
\code{0x100104fc} & \code{player\_sprite} & 9 words & the $3 \times 3$ character sprite\\
\bottomrule
\end{tabular}
\caption{Layout of the \code{.data} section (it ends at \code{0x10010520}).}
\label{tab:memory}
\end{table}

\subsection{Sprites and pixel grids}

\textbf{Bubbles.} A bubble is drawn from one entry of \code{palette}, which
holds a dark shade $D$, the base colour $B$ and a light shade $L$:
\begin{center}
\begin{tabular}{|c|c|c|}
\hline $D$ & $B$ & $D$\\
\hline $B$ & $L$ & $B$\\
\hline $D$ & $B$ & $D$\\
\hline
\end{tabular}
\end{center}
Entries 1 to 4 are the four bubble colours. Entry 0 is all black, so drawing
``colour 0'' erases a cell; this is why the whole field can simply be redrawn
every frame.

\textbf{Character.} \code{player\_sprite} stores the nine colours of the
character in row major order. Its centre unit is replaced by the colour of
the held bubbles while the character is holding something.

\subsection{Memory after assembling}

\begin{figure}[H]
\centering
\figurefile{\textwidth}{figures/memory_screenshot.png}
\caption{Memory in Saturn right after assembling (the program is stopped at a
breakpoint on the first instruction of \code{main}). Top: from
\code{0x10010000}, the two addresses followed by the empty \code{grid}.
Bottom: from \code{0x100104a0}, the variables (\code{player\_col} $= 3$,
\code{row\_interval} $=$ \code{0x1f4} $= 500$), the \code{palette} from
\code{0x100104c0} and \code{player\_sprite} from \code{0x100104fc}.}
\label{fig:memory}
\end{figure}

The values to look for in the screenshot are listed in
Table~\ref{tab:expected}.

\begin{table}[H]
\centering
\small
\begin{tabular}{@{}lll@{}}
\toprule
Address & Value & Meaning\\
\midrule
\code{0x10010000} & \code{0x10008000} & \code{display\_address}\\
\code{0x10010004} & \code{0xffff0000} & \code{keyboard\_address}\\
\code{0x10010008} -- \code{0x100104a0} & \code{0x00000000} & \code{grid}, \code{mark}, \code{ff\_stack}, \code{ff\_count}\\
\code{0x100104a4} & \code{0x00000003} & \code{player\_col} (the middle column)\\
\code{0x100104b8} & \code{0x000001f4} & \code{row\_interval} $= 500$\\
\code{0x100104bc} & \code{0x00000000} & \code{frame\_count}\\
\code{0x100104cc} & \code{0x00900000} & \code{palette} entry 1 (red), dark shade\\
\code{0x100104d0} & \code{0x00e02020} & \code{palette} entry 1 (red), base colour\\
\code{0x100104d4} & \code{0x00ff9090} & \code{palette} entry 1 (red), light shade\\
\code{0x1001050c} & \code{0x00ff60c0} & centre of \code{player\_sprite}\\
\bottomrule
\end{tabular}
\caption{Values expected in memory right after assembling.}
\label{tab:expected}
\end{table}

\section{Milestone 1: The Static Scene}

\code{draw\_walls} clears the screen and draws what never changes: the four
walls and the red line above the character's row. \code{draw\_board} draws
what can change: the 98 cells of the field (\code{draw\_grid}), the aiming
line (\code{draw\_aim}) and the character (\code{draw\_player}). They are built on \code{draw\_bubble}
($x$, $y$, colour), \code{draw\_hline}, \code{draw\_vline} and
\code{clear\_screen}.

A game starts with five rows of random bubbles in four colours: \code{main}
calls \code{push\_row} five times. Random rows can already contain a vertical
match-3, so \code{main} then runs \code{resolve\_matches} once; this is why
some columns may start shorter than five bubbles.

\begin{figure}[H]
\centering
\figurefile{0.3\textwidth}{figures/static_scene.png}
\caption{Screenshot of the Bitmap Display in Saturn at the start of a game:
the $7 \times 14$ playing field with five rows of random bubbles (one column
lost bubbles to a match in the random rows), the flashing aiming line and the
character in the bottom row.}
\label{fig:static}
\end{figure}

\section{Milestone 2: Movement, Grabbing and Row Advancement}

The game loop follows the order suggested in the handout: check the keyboard,
process the action, check for row advancement, redraw, sleep for 16\,ms
(about 60 frames per second) and repeat. The keyboard is polled at
\code{0xffff0000}; the key is read from \code{0xffff0004}. Upper case letters
are accepted as well.

\begin{itemize}[nosep]
\item \textbf{Movement.} \code{a} and \code{d} change \code{player\_col}.
\item \textbf{Grabbing.} \code{s} calls \code{handle\_grab}. It finds the
lowest bubble of the character's column (\code{find\_lowest}) and removes it
together with all bubbles of the same colour stacked directly above it. The
colour and the count are stored in \code{held\_colour} and
\code{held\_count}. Pressing \code{s} again adds more bubbles, but only if
they have the colour that is already held.
\item \textbf{Row advancement.} Every \code{a}, \code{d}, \code{s} or
\code{w} increases \code{key\_count}; after 10 keystrokes a new row arrives.
\code{push\_row} moves every bubble down by one row and fills row 0 with
seven random bubbles (system call 42).
\item \textbf{Repainting.} \code{draw\_board} repaints the field every frame.
\item \textbf{Quit.} \code{q} ends the program with system call 10.
\item \textbf{Game over.} After a new row has been placed,
\code{advance\_row} checks row 13. If it contains a bubble, \code{game\_over}
draws the field one last time and the program ends. In this last frame the
character is only drawn if its own cell is still empty, so a bubble that has
reached the character's row is never hidden behind the character.
\end{itemize}

\section{Milestone 3: Throwing, Matching and Chain Reactions}
\label{sec:m3}

\textbf{Position limits.} \code{move\_left} and \code{move\_right} do nothing
when the character is already in column 0 or column 6.

\textbf{Throwing.} \code{w} calls \code{handle\_throw}. The landing position
is the first free row of the character's column, directly below its lowest
bubble. The held bubbles are placed from there downwards. They may use rows 0
to 12 only; bubbles that do not fit stay in the character's hand, and if
there is no free row nothing happens.

\textbf{Matching (hard features 2 and 3).} After every throw, and after
every new row, the function \code{resolve\_matches} runs the following
loop:

\begin{enumerate}[nosep]
\item \textbf{Trigger} (\code{find\_matches}). Every column is scanned from
top to bottom for runs of three or more bubbles of the same colour. The cells
of such a run are marked in \code{mark} and pushed on \code{ff\_stack}. If
nothing was marked, the loop ends.
\item \textbf{Explosion: flood fill} (\code{flood\_fill}, hard 2). While the
stack is not empty, a cell is popped and its four neighbours (above, below,
left, right) are examined. A neighbour that has the same colour and is not
marked yet is marked and pushed. Every cell is pushed at most once, so the
stack never needs more than 98 entries.
\item The marked bubbles are removed (\code{remove\_marked}).
\item \textbf{Gravity} (\code{float\_step}, hard 3). Every bubble that has an
empty cell directly above it moves up by one row. This step is repeated until
nothing moves any more, so every column is again filled from the top without
gaps.
\item The loop starts again at step 1: if gravity has lined up a new
vertical match-3, it is removed as well.
\end{enumerate}

Figure~\ref{fig:screens} shows the loop on an example. The steps happen
within one frame; the pictures b) to e) were drawn by the game's own drawing
code after each step.

\begin{figure}[H]
\centering
\figurefile{0.8\textwidth}{figures/screens.png}
\caption{a) The character holds a blue bubble under column 2; the dotted
aiming line shows where it will land. b) The throw completes a vertical blue
match-3. c) The flood fill also takes the blue bubble connected to it in
column 3, so four blues are removed. d) Gravity lets the two reds in column 3
float up, lining up three reds. e) The loop runs again and removes them. f) A
game that is over: a new row has pushed bubbles into the character's row;
the bubble in the character's cell is shown instead of the character.}
\label{fig:screens}
\end{figure}

\section{Milestones 4 and 5: Features}

\textbf{Timer (easy 1).} Every frame without a new row increases the counter
\code{row\_timer}. When it reaches the value of \code{row\_interval}, a new
row arrives. A new row from the timer or from 10 keystrokes resets both
counters, so the rows come from whichever happens first.

\textbf{Faster and faster (easy 2).} After every new row, \code{row\_interval}
becomes 50 frames shorter: 500, 450, 400, \dots, down to 150 frames. At about
60 frames per second this is roughly 8 seconds at the start and 2.5 seconds
at the end.

\textbf{Aiming line (hard 6).} \code{draw\_aim} draws one grey dot in the
centre of every empty cell of the character's column, from the lowest bubble
down to the red line, so it shows where a throw will land. It only draws when
bit 4 of \code{frame\_count} is 0, which is true for 16 frames and false for
the next 16, so the line flashes. \code{draw\_grid} paints the empty cells
black first, so the dots disappear by themselves in the frames where they are
not drawn.

\textbf{Flood fill and gravity (hard 2, 3).} See Section~\ref{sec:m3}.

\section{How to Play}

\textbf{Setup.} Open \code{MagicalDrop.asm} in Saturn or MARS. Set the Bitmap
Display to a unit width and height of 8, a display width of 256, a display
height of 512 and the base address \code{0x10008000 (\$gp)}. In MARS, connect
both the Bitmap Display and the Keyboard and Display MMIO Simulator to MIPS.
Assemble and run the program, then click the bitmap window (Saturn) or the
keyboard area (MARS) so that it receives the keystrokes.

\textbf{Controls.}

\begin{center}
\begin{tabular}{@{}cl@{}}
\toprule
Key & Action\\
\midrule
\code{a} / \code{d} & move the character left / right\\
\code{s} & grab the lowest bubble of the column above, with all bubbles of the same\\
 & colour directly above it; press again to add more bubbles of that colour\\
\code{w} & throw the held bubbles up into the column above (where the dotted line ends)\\
\code{q} & quit\\
\bottomrule
\end{tabular}
\end{center}

\textbf{Goal.} Line up three or more bubbles of the same colour vertically by
throwing bubbles into a column. The match removes those bubbles and every
bubble of the same colour connected to them. The bubbles below float up and
can line up new matches. A new row pushes the field down after every 10
keystrokes or when the row timer runs out, whichever comes first, and the
timer gets shorter after every new row. The game is over when a bubble
reaches the bottom row, where the character moves.

\textbf{Screen.} The centre of the character shows the colour of the bubbles
it is holding. The flashing dotted line above the character shows where a
throw will land. The red line marks the character's row: a bubble below it
ends the game.

\section{Attribution Table}
\label{sec:attribution}

Table~\ref{tab:attribution} lists which team member is responsible for each
part of \code{MagicalDrop.asm}, and Table~\ref{tab:attribution-report} the
parts of this report and of the demonstration. Each milestone has one person
responsible; Xiangyu Zhou also connects both parts in \code{main} and
\code{game\_loop}.

\begin{table}[H]
\centering
\small
\begin{tabular}{@{}>{\raggedright\arraybackslash}p{2.0cm}>{\raggedright\arraybackslash}p{5.0cm}>{\raggedright\arraybackslash}p{5.6cm}>{\raggedright\arraybackslash}p{2.1cm}@{}}
\toprule
Part & Work & Code & Student\\
\midrule
Milestone 1 & Static scene: walls, red line, grid, bubbles, character and
colours & \code{draw\_walls}, \code{draw\_board}, \code{draw\_grid},
\code{draw\_bubble}, \code{draw\_player}, \code{draw\_hline},
\code{draw\_vline}, \code{clear\_screen}, \code{palette},
\code{player\_sprite} & Xiangyu Zhou\\
\addlinespace
Milestone 2 & Game loop, keyboard, movement, grabbing, the five starting rows,
a new row every 10 keystrokes, repainting, quitting, game over &
\code{main}, \code{game\_loop}, \code{move\_left}, \code{move\_right},
\code{handle\_grab}, \code{find\_lowest}, \code{push\_row},
\code{random\_colour}, \code{advance\_row}, \code{game\_over},
\code{quit\_game} & Xiangyu Zhou\\
\addlinespace
Milestone 3 & Position limits, throwing, finding and removing vertical
match-3s & the limits in \code{move\_left} and \code{move\_right},
\code{handle\_throw}, \code{find\_matches}, \code{resolve\_matches},
\code{remove\_marked} & Belina Zhu\\
\addlinespace
Easy 1 & Row advancement on a timer & \code{row\_timer} in \code{game\_loop}
& Xiangyu Zhou\\
\addlinespace
Easy 2 & Row advancement gets faster & the lines in \code{game\_loop} that
shorten \code{row\_interval} & Belina Zhu\\
\addlinespace
Hard 2 & Flood fill & \code{flood\_fill}, \code{ff\_try}, \code{ff\_stack} &
Belina Zhu\\
\addlinespace
Hard 3 & Gravity: unsupported bubbles float up & \code{float\_step} &
Belina Zhu\\
\addlinespace
Hard 6 & Flashing dotted aiming line & \code{draw\_aim},
\code{frame\_count} & Xiangyu Zhou\\
\addlinespace
Integration & Connecting both parts, the \code{.data} layout, the preamble,
testing in Saturn & \code{main}, \code{game\_loop}, \code{.data} &
Xiangyu Zhou\\
\bottomrule
\end{tabular}
\caption{Who is responsible for each part of \code{MagicalDrop.asm}.}
\label{tab:attribution}
\end{table}

\begin{table}[H]
\centering
\small
\begin{tabular}{@{}>{\raggedright\arraybackslash}p{2.1cm}>{\raggedright\arraybackslash}p{7.0cm}>{\raggedright\arraybackslash}p{5.6cm}@{}}
\toprule
Student & Report sections & Demonstration\\
\midrule
Xiangyu Zhou & Bitmap Display Configuration, Memory Layout Plan (with the
memory screenshots), Milestone 1, Milestone 2, easy 1, hard 6, How to Play &
starting screen, movement, grabbing, rows after 10 keystrokes, timer rows,
aiming line, game over, quitting\\
\addlinespace
Belina Zhu & Milestone 3, easy 2, hard 2, hard 3 (with the step-by-step
figure) & position limits, throwing, removing a vertical match-3, flood fill,
bubbles floating up, rows coming faster\\
\bottomrule
\end{tabular}
\caption{Who is responsible for each part of this report and of the
demonstration.}
\label{tab:attribution-report}
\end{table}

\end{document}
